\par\SourceLeaf{152}
Es bleibt mir jetzt noch übrig, einige zum Theil schon oben berührte Punkte noch etwas näher zu betrachten.

Man hat oben gesehen, wie über ein Grundprincip P. und R. ganz entgegengesetzte Ansichten gehabt haben (S. [133] und [137]. Wenn der letztere S. [138, oben] von einem Steigen und \textit{Sinken} des \textit{Fonds} spricht, so vermuthe ich, dass er eigentlich nur den \textit{Ertrag} des Fonds gemeint hat. Denn das scheint mir, insofern die Anstalt ein Beneficium ist, die strenge Pflicht der Verwaltung zu sein, dafür zu sorgen, dass die Substanz, aus welcher das Beneficium fliesst, in ihrer Integrität erhalten werde. Dies kann aber schon wegen der bei Kapitalausleihungen von Zeit zu Zeit bei aller Vorsicht nicht abzuwendenden Verluste mit Sicherheit anders nicht geschehen, als wenn man neben der Erhaltung auch einige allmählige Vermehrung sich zum Ziele setzt, wobei man denn immer lieber etwas zu viel als zu wenig thun möge. Auf diese Weise wird die \textit{Gesammtheit} der Percipienten in der spätern Zeit gegen die Gesammtheit der frühern nicht zu kurz kommen, sondern vielmehr eher besser daran sein, was aber die dermaligen Percipienten jenen um so eher gönnen können, da sie selbst die Früchte einer ähnlichen Enthaltsamkeit \textit{ihrer} Vorgänger geniessen. Ausserdem erfordert die Billigkeit, dass man sich bestrebe, das unvermeidliche und bei einer kleinen Gesellschaft verhältnissmässig sehr grosse Schwanken der Zahl der Percipienten durch zweckmässige Maassregeln vo viel thunlich auszugleichen. Dagegen aber scheint mir P.’s Forderung, dass niemals ein künftiger \textit{einzelner} Percipient weniger erhalten solle, als irgend ein früherer erhalten hat, bei einer Gesellschaft, die wie \textsc{Sch.} sehr richtig bemerkt hat, immer etwas Actienmässiges behalten wird, im Rechte nicht begründet; jedenfalls aber, und diess ist der Hauptpunkt auf den es ankommt, lässt sich einer solchen Forderung, wenn sie wie eine unbedingte gelten soll, gar nicht genügen ohne die offenbarste Unbilligkeit gegen die dermaligen Percipienten. Es liegt auf der Hand: je grössere Sicherheit man verlangt, dass jener Fall niemals eintreten müsse, desto weniger darf man den jetzigen Percipienten verabreichen. Man müsste die extremsten Fälle für die mögliche Zahl der Percipienten berücksichtigen, wovon, wie oben gezeigt ist, P.’s Ansätze weit entfernt waren. Noch viel schlagender tritt dies hervor durch die weiter unten S. [150] aufgestellten Überschläge auf den Grund der erweiterten Interessentenzahl, einer Eventualität, die doch auch schon 1794 unter die Zahl der künftig nicht bloss möglichen, sondern sogar wahrscheinlichen Fälle hätte aufgenommen werden können. Denn damals war die Zahl aller Universitätsprofessoren 45\footnote{Einer davon, K., starb noch vor Beendigung der Verhandlungen 1794 August 21. Jetzt (im Herbst 1845) sind 57. Aber die Anzahl der verheiratheten Mitglieder der Witwenkasse ist in viel stärkerm Verhältniss gestiegen von 26 im Jahr 1794 auf 42 im Jahr 1845.}, funfzig Jahre früher nur 25, und welche Grenzen die fortwährend gesteigerten Zeitbedürfnisse finden werden, ist unmöglich im Voraus festzusetzen.

Der zweite Punkt betrifft die Auslegung der Progressionsnormirung in \textsc{Sch.}’s Sinn. Wenn \textsc{Sch.} sagt, [S. 138. Z. 19] P. habe das Maass der Erhöhung und eben so der Verminderung durch eine unwandelbare Regel bestimmt, so setzt diess zwar nothwendig die erste Interpretation S. [129] voraus, aber doch räumt \textsc{Sch.} damit zu viel ein. Eine unwandelbare Regel für die Verminderung fand sich darin nur in so weit, als bestimmt wurde, \textit{wann} Verminderung eintreten müsse (nemlich, nach jener Auslegung, sofort nachdem die Zahl von 15 Pensionen überschritten), aber noch nicht für die \textit{Grösse} der Verminderung. Diese Unvollständigkeit hätte, deucht mir, nach damaliger Lage der Sache, am füglichsten durch eine Bestimmung in folgender Fassung ergänzt werden können:
\begin{quote}
So oft das Kapitalvermögen um 5000 Rthl. gestiegen ist, tritt eine Erhöhung der Pensionen ein, welche für jede einzelne Pension 10 Rthl. beträgt, wenn und so lange nicht mehr als 15 Pen%
